\documentclass[article,A4,12pt]{llncs}
\usepackage[T1]{fontenc}
\usepackage{amsmath}
\usepackage{amssymb}
\usepackage{amsfonts}
\usepackage{mathrsfs, bm}

\usepackage{graphicx}
\usepackage{tabularx}
\usepackage{subfig}
\usepackage{epsf,times}
\usepackage{color}
\usepackage{wrapfig}
\usepackage{cases}
\usepackage{multicol}

\usepackage[T1]{fontenc}
%\newcommand{\tmname}[1]{\textsc{#1}}
%\newcommand{\tmop}[1]{\ensuremath{\operatorname{#1}}}
%\newcommand{\tmsamp}[1]{\textsf{#1}}
%\newcommand{\tmtextsc}[1]{{\scshape{#1}}}
%\newcommand{\tmtextsl}[1]{{\slshape{#1}}}
%\newcommand{\tmtexttt}[1]{{\ttfamily{#1}}}

\leftmargin=0.0cm
\oddsidemargin=0.5cm
\evensidemargin=0.5cm
\topmargin=0cm
\textwidth=16.0cm
%\textheight=21.5cm
\textheight=20.0cm
\pagestyle{plain}
\setlength{\columnsep}{20pt}

\def\m{\mathbf{m}}
\def\H{\mathbf{H}}
\def\E{\mathbf{E}}
\newcommand{\vepsi}{{\varepsilon}}
\def\hnorm#1#2{\vert\,#1\,\vert_{#2}}
\newcommand{\R}{{\mathbb R}}
\newcommand{\Sph}{{\mathbb S}}
\def\x{\mathbf{x}}
\def\hvec{\overline{\mathbf{h}}}
\def\evec{\overline{\mathbf{e}}}

\newcommand{ \etal}{\mbox{\emph{et al. }}}

\newcommand\vect[1]{\mbf{#1}}
\newcommand{\mbf}[1]{\mbox{\boldmath$#1$}} 
\newcommand{\RC}[1]{#1 $\times$ #1 $\times$ #1}
\def\um{$\mu$m}
\def\C{$^{\circ}\mathrm{C}$}

\newcommand{\Rmnum}[1]{\expandafter\@slowromancap\romannumeral #1@}

% DEFINITION OF CUSTOM FONT SIZE
\newcommand{\customfontA}{\fontsize{50}{55}\selectfont}
\newcommand{\customfontB}{\fontsize{14.4}{20}\selectfont}
\newcommand{\customfontC}{\fontsize{30}{35}\selectfont}

\DeclareMathAlphabet{\mathpzc}{OT1}{pzc}{m}{it}

\def\clovek#1{\noindent\bgroup\vbox{\noindent#1}\egroup\vskip1em}

% TO INPUT BACKGROUND IMAGE
\usepackage{eso-pic}
\newcommand\BackgroundPic{
\put(0,0){
\parbox[b][\paperheight]{\paperwidth}{
\vfill
\centering
\includegraphics[width=\paperwidth,height=\paperheight]{img/intro-frontpage.png}
%\includegraphics[width=\paperwidth,height=\paperheight]{img/background.jpg}
\vfill
}}}

\begin{document}

% INPUTTING BACKGROUND IMAGE
\AddToShipoutPicture{\BackgroundPic}
\vbox{}
\pagestyle{empty}
\newpage
\textwidth=15.5cm
\ClearShipoutPicture
\newpage

%%%%%%%%%%%%%%%%%%%%%%%%%%%%%%%%%%%%%%%%%%%%%%%%%%%%%%%%%%%%%%%%%%%%%%%%%

\section*{}
\small
\input ../common/aboutnclab.tex

\subsection*{Acknowledgement}
This publication was created with the help of numerous freely 
available web resources and tutorials related to Python, Scipy,
Numpy, Pylab, Matplotlib, Sympy and other projects.

\normalsize

\newpage
%{\ }
\setcounter{tocdepth}{2}
\tableofcontents
%\pagestyle{plain}

\newpage

\pagestyle{plain}
\setcounter{page}{1}

%%%%%%%%%%%%%%%%%%%%%%%%%%%%%%%%%%%%%%%%%%%%%%%%%%%%%%%%%%%%%%%%%%%%%%%%%

\section{What Is Cloud Computing?}

Cloud Computing sounds like a rocket science, but it is not,
and it really changes things for the better. Chances are that you are using 
it already, without even noticing it. For example, if you are using Gmail, 
Hotmail, Yahoo or another such provider to do your emails, then your data are 
stored on the cloud. 

\subsection{Some Trivia}

In the IT sense of the word, and with a bit of simplification, the {\em Cloud} 
is a pool of interconnected computers. They come in different sizes depending on the provider - 
large cloud facilities have millions of computing cores but some companies are running 
their own private clouds with relatively few computers. The computers in large clouds
are not very similar to your desktop PC -- they are stripped off many unnecessary 
things, compacted, and interconnected into a powerful grid. Their processors 
typically contain multiple {\em computing cores} that can share the same memory.
Such a {\em multicore processor} is much more efficient compared to a cluster 
of single-core PCs with the same number of cores.

The main advantage of the cloud is its {\em elasticity}. Upon your request, the
provider will assemble for you an {\em instance} with the parameters that 
you need. An instance is created within a minute or so, and the user can upgrade 
or downgrade it dynamically, depending on the actual needs. After the instance is 
no longer needed, it is dissolved and the same hardware is used to build instances 
for other users. The following is a realistic example of how such instances may 
look like (just the name of the provider was omitted):\\

\begin{center}
\begin{tabular}{|l|l|l|l|}
\hline
{\bf Small instance} & 1.7 GB of memory & 1 core & 160 GB hard disk \\
\hline
{\bf Medium instance} & 7.5 GB of memory & 4 cores & 850 GB hard disk \\
\hline
{\bf Large instance} & 15 GB of memory & 8 cores & 1690 GB hard disk \\
\hline
\end{tabular}
\end{center}

\vspace{4mm}
\noindent
Importantly - these three instances are created by grouping the same resources together 
in different ways. 
It is quite interesting to look at the prices:

\begin{center}
\begin{tabular}{|l|l|}
\hline
{\bf Small instance} &	\$0.085 per hour\\
\hline
{\bf Medium instance}&	\$0.34 per hour	\\
\hline
{\bf Large instance}&	\$0.68 per hour\\
\hline
\end{tabular}
\end{center}

\vspace{4mm}
\noindent
With such low prices, literally anyone has now access to 
a powerful computer -- that otherwise would cost many thousands of dollars -- 
for just cents per hour. This brings new wonderful opportunities to 
education at all levels.

\subsection{Software as a Service (SaaS)}

Another big change that the Cloud Computing brings is in how software is handled. 
Traditionally, one would buy a software in a big box with a small CD in it, 
and install it on one's computer. This model is now being challenged by a new 
approach called {\em Software as a Service (SaaS)}. The user does not have 
to own a copy of the software physically. Instead, the software is running 
on a remote server and accessed by users over the Internet. Typically, paying 
for the access is much less expensive compared to buying the software.

\subsection{Access from Mobile Platforms}

Last but not least, since the software is running on the cloud, the user's hardware 
does not matter so much anymore. The only really important things is to be able to 
run a web browser and access the Internet. This can be done on anything between smart 
phones, tablets, netbooks, laptops, and desktop computers. 

\subsection{Some Myths}

Very often, Cloud Computing is mentioned in the context of large 
computations in science, engineering, finance, and other fields. This is because there are many 
cloud providers and they need to show off. But in reality, the clouds are mostly
busy processing small tasks of ordinary users. In that case, you may say, a standard 
office PC should be enough. Not exactly. Once you are not in your office, for example
while traveling, it is difficult to reach your office PC and do things as usual. 
The cloud changes that -- it allows you to access your data 
from everywhere, any time. Once you get used to it, there is no way back. 

\subsection{NCLab and K-12 Education}

NCLab is a web-based framework that provides 
instant access to many programming, mathematics, computer
modeling, and scientific computing activities. NCLab does not have to be installed - it 
is automatically available in any classroom that has Internet access. In 
contrast to traditional educational software products, 
users can access their accounts and work from anywhere and at any time.
Students can start a problem at school, finish the rest at home, and notify 
the instructor about the completion with one mouse click. Most importantly, 
however, NCLab invites K-12 teachers and students to discover new exciting 
areas of mathematics, programming, and computer modeling that were not 
available to them before.

%%%%%%%%%%%%%%%%%%%%%%%%%%%%%%%%%%%%%%%%%%%%%%%%%%%%%%%%%%%%%%%%%%%%%%%%%
\newpage

\section{Getting Started}

In order to explore NCLab, first visit its home page {\tt http://nclab.com}.
Under the logo you will see a slide show representing selected activities 
offered by NCLab. The Cloud Monitor, located under the slide show, displays 
the current load on the core cluster. By default, NCLab runs on six nodes 
with a total of 100 computing cores, and has approximately 100 GB of memory 
available for immediate use. These parameters can be scaled up dynamically 
if needed. 

The User Statistics window, located in the 
bottom-right corner, provides current overview of NCLab's users and their projects. 
The window above User Statistics is a WebGL tester. WebGL is a very recent 
web-browser technology that supports amazing 3D graphics. If you can see in your 
browser an image similar to the one in Fig. \ref{fig:outside}, then your web 
browser supports it. In that case, click into the WebGL tester window, and move 
your mouse while keeping the left button pressed, to see WebGL in action.
If your browser does not support WebGL, we recommend that you 
upgrade it to a newer version. IE is one of few browsers which do not support WebGL. Therefore 
it is better to use Chrome, Firefox, Opera, or Safari. A login dialog where
users enter their username, password and institution code is located in the 
top-right corner.\\[-8mm]

\begin{figure}[!ht]
\begin{center}
\includegraphics[width=\textwidth]{img/outside.png}
\end{center}
\vspace{-2mm}
\caption{NCLab's home page.}
\label{fig:outside}
\vspace{-0.6cm}
\end{figure}

\newpage

\subsection{Creating an Account}

In order to create a new account, click on the link "Create free account" which is next to
"Username". The following window appears:

\begin{figure}[!ht]
\begin{center}
\includegraphics[width=0.7\textwidth]{img/create-account.png}
\end{center}
\vspace{-4mm}
\caption{Creating new account.}
\label{fig:creacc}
\end{figure}

\noindent
In this form, enter your preferred username, email address (to be used only if you ask
us to generate a new password for you), and password twice. Next:

\begin{itemize}
\item If you belong to an institution 
      that uses NCLab for teaching or research, you should know your institution's code. Enter it in the form, 
      along with your first and last names. Your name will be visible to an NCLab admin
      who will verify that you indeed belong to the institution.
\item If you are using NCLab for your personal hobby or self-education, for activities not 
      related to any institutional use, then leave the last three lines empty. 
\end{itemize}
After clicking OK, you are ready to login!

\subsection{After Login}

The interior of NCLab is similar to the standard computer desktop. Many users maximize 
the web browser window to make the desktop experience even more realistic.

If any member of any of your groups is in NCLab at the time you log in, the Groups icon 
will be lighted green as shown in Fig. \ref{fig:desktop}. 

\begin{figure}[!ht]
\begin{center}
\includegraphics[width=\textwidth]{img/desktop.png}
\end{center}
%\vspace{-2mm}
\caption{NCLab desktop after login.}
\label{fig:desktop}
\end{figure}

\noindent
More about Groups and Chat will be said in a moment. First, let us briefly 
describe the other icons:

\begin{itemize}
\item File Manager is used to manage files and folders similarly to Windows Explorer. 
\item Programming Module enables programming in Karel (popular educational language for beginners),
      Python (modern programming language for advanced applications), and Javascript 
      (most popular language for web development). Additional languages are in preparation.
\item Fractal Explorer is an interactive graphical application that allows the users
      to learn about Mandelbrot and Julia fractals, and to generate beautiful fractal art.
\item Settings: Here the user can upload his/her avatar image (to be used in Groups and Chat), 
      request a new password, and adjust sound preferences.
\item Help launches a quick help about NCLab.
\item Logoff will close all applications and end the session.
\end{itemize}

\subsection{Groups and Chat}

In NCLab, users can form groups, see who is online, and chat in real time. In order to 
create a new group, click on the Groups icon. A window similar to the one shown 
in Fig. \ref{fig:groups} appears. This concrete screenshot shows four groups of the 
user "solin" -- the top one was created by Jordan, the other three by solin. The 
green light next to the last group indicates that someone in that group
is currently in NCLab. Before we go look who that is, note that you can create a
new group of your own by clicking on "New Group". In order to add people to your groups,
you need to know their usernames. The button "Who's online" allows you to see instantly 
all users in all your groups who are currently logged in. 

\begin{figure}[!ht]
\begin{center}
\includegraphics[width=0.5\textwidth]{img/groups.png}
\end{center}
%\vspace{-2mm}
\caption{Groups window, showing that someone in the last group is currently in NCLab.}
\label{fig:groups}
\end{figure}

\newpage
\noindent
Clicking on the last group 
reveals that the online users are Emily and Rick.
\begin{figure}[!ht]
\begin{center}
\includegraphics[width=0.5\textwidth]{img/groups-2.png}
\end{center}
%\vspace{-2mm}
\caption{Detail of the last group.}
\label{fig:groups-2}
\end{figure}

\noindent
Clicking on any user who has the green light next to him/her (except yourself),
you initiate a chat to that user. Or someone from your groups may contact you. 
The Chat window is shown in Fig. \ref{fig:chat}.

\begin{figure}[!ht]
\begin{center}
\includegraphics[width=0.4\textwidth]{img/chat.png}
\end{center}
%\vspace{-2mm}
\caption{Chat window.}
\label{fig:chat}
\end{figure}


\end{document}
